% !TEX root = ../main.tex
%% ---------------------------------------------------------------------------
%% 3. Error estimate of Galerkin meshfree formulation
%% ---------------------------------------------------------------------------
\section{Error estimate of Galerkin meshfree formulation}\label{sec:error}


\subsection{$H_2$ norm of standard error for Galerkin meshfree formulation}\label{sec:h2standard}

In \cite{ref50}, $\|u-u^{h}\|_{H_2}$ norm of Galerkin formulation, especially for Galerkin finite element methods analysis has been investigated.
However, since the integration error of Galerkin meshfree formulation can not be ignored, the conclusion in \cite{ref50} become not fit here.
More specifically, the $\|u-u^{h}\|$ norm of Galerkin meshfree formulation for Kirchhoff--Love plate theory can be stated as \cite{ref50}:

In order to estimate the standard error $e$ of Galerkin weak form \eqref{weakform}, combining Eq. \eqref{nitsche_weakform} and Eq. \eqref{weakform} gives the orthogonality condition:
\begin{equation}\label{orthogonality_std}
\mathrm{A}(v^{h},w-w^{h})=0,\quad \forall v^{h}\in V^{h}
\end{equation}
Invoking continuity \eqref{continuity}, coercivity \eqref{coercivity} and orthogonality \eqref{orthogonality_std} conditions, we obtain the $H_2$ norm of standard error $e$:
\begin{equation}\label{h2_standard}
\begin{aligned}
\alpha \Vert e \Vert_{H_2}^2 &\le \vert a(e,e) \vert \\
&\le \left \vert a(e,w-w^h) + \underbrace{a(e,w^h-\mathcal I^p w) + Ch}_{\textrm{Orthogonality of Eq. \eqref{orthogonality_std}}} \right \vert \\
&\le \vert a(e,w-\mathcal I^p w) \vert + Ch \Vert e \Vert_{H_2} \\
&\le Ch \Vert e \Vert_{H_2} \Vert w-\mathcal I^p w \Vert_{H_2} + Ch \Vert e \Vert_{H_2} \\
&\le (Ch^p + Ch^q) \Vert e \Vert_{H_2}
\end{aligned}
\end{equation}
equivalently
\begin{equation}
	\Vert e \Vert_{H_2} \le Ch^p + Ch^q
\end{equation}

Suppose that the following second order problem is introduced herein as the duality problem of Nitsche trick.
\begin{equation}
	a(v,u) = (v,e), \quad \forall v \in H_2
\end{equation}
while $\Delta^2 u = e$ in $\Omega$, and $e = w-w_h \in H_2$.
According to closed graph theorem, we have:
\begin{equation}
	\vert u \vert_{H_k} \le C \Vert e \Vert_{L_2}, \quad \forall k \le 4
\end{equation}

\begin{equation}
	\begin{aligned}
		\Vert e \Vert_{L_2}^2 &= (e,e)_{L_2} = a(e,u) & \\
		&= a(e,u - \mathcal I^p u) & (\textrm{Orthogonality of \eqref{orthogonality_std}}) \\
		&\le C \Vert e \Vert_{H_2} \Vert u - \mathcal I^p u \Vert_{H_2} & (\textrm{Continuousity of \eqref{continuity}})
	\end{aligned}
\end{equation}

\begin{equation}
	\Vert u - \mathcal I^p u \Vert_{H_2} \le C h^{l-1} \vert u \vert_{H_{l+1}}, \quad \forall l \le p
\end{equation}


As $p=2$,
\begin{equation}
	\Vert e \Vert_{L_2}^2 = C h \Vert e \Vert_{H_2} \vert u \vert_3
	\le C h \Vert e \Vert_{H_2} \Vert e \Vert_{L_2} 
	\Rightarrow \Vert e \Vert_{L_2} = C h \Vert e \Vert_{H_2}
\end{equation}
As $p>2$,
\begin{equation}
	\Vert e \Vert_{L_2}^2 = C h^2 \Vert e \Vert_{H_2} \vert u \vert_4
	\le C h^2 \Vert e \Vert_{H_2} \Vert e \Vert_{L_2} 
	\Rightarrow \Vert e \Vert_{L_2} = C h^2 \Vert e \Vert_{H_2}
\end{equation}
% Further introducing error estimate of meshfree approximation into Eq. \eqref{h2_standard} naturally gives:
% \begin{equation}\label{eq:49}
% \|e\|_{H_2}\le C_{1p}h^{k-1}|w|_{H_{k+1}},\quad k=\min\{p,m\}
% \end{equation}
% where $h$ denotes the characterized length of support domain, $C_{1p}$ is a constant related to $\alpha$ and $\gamma$.

% The estimate of Eq. \eqref{eq:49} is consistent with \cite{ref46}.
% This is because the same property of Galerkin weak form \eqref{eq:26}.
% However, in real computation, the numerically discrete Galerkin meshfree weak form \eqref{eq:35} is adopted, which leads to $\eta\neq0$.
% Subsequently, we'll assess $\|\eta\|_{H_2}$ through orthogonality condition accounting for numerical integration.

% \subsection{$H_2$ norm of error for Galerkin meshfree formulation to Kirchhoff--Love plate}\label{sec:h2error}

% With consideration for the effect of numerical integration, the orthogonality condition will no longer holds for $\tilde{e}$:
% \begin{equation}\label{eq:50}
% \mathrm{A}(v^{h},\tilde{e})\neq0,\quad \forall v^{h}\in V^{h}
% \end{equation}
% Integration constraint plays a crucial role in numerical integration of Galerkin meshfree methods.
% Following Eq. \eqref{eq:37}, the $q$th order integration constraint can be rewritten by:
% \begin{equation}\label{eq:51}
% \mathbf{A}^{*}(v^{h},T^{q}w)=\mathbf{B}^{*}(v^{h},T^{q}w),\quad \forall v^{h}\in V^{h}
% \end{equation}
% where $T^{q}w\in\mathbb{P}^{q}$, $\mathbb{P}^{q}=\operatorname{span}\{\phi_{ijk}\}_{i+j+k=0}^{q}$, $q\le p$ denotes the $q$th order Lagrange interpolant of $w$, whose corresponding interpolation error is bound by \cite{ref51}:
% \begin{equation}\label{eq:52}
% \left\|w-T^{q}w\right\|_{H_l}\le Ch^{k-l+1}|w|_{H_{k+1}},\quad l\le k+1,\;k=\min\{q,m\}
% \end{equation}
% And, $\mathbf{A}^{*}(v^{h},T^{q}w)$ and $\mathbf{B}^{*}(v^{h},T^{q}w)$ in Eq. \eqref{eq:51} are listed as follows:
% \begin{equation}\label{eq:53}
% \mathbf{A}^{*}(v^{h},T^{q}w)=\int_{\Omega}^{*}B^{T}(v^{h})M(T^{q}w)\mathrm{d}\Omega
% \end{equation}
% \begin{equation}\label{eq:54}
% \begin{aligned}
% \mathbf{B}^{*}(v^{h},T^{q}w)={}&-\int_{\Gamma}v^{h}V(T^{q}w)\mathrm{d}\Gamma+\int_{\Omega}v^{h}q(T^{q}w)\mathrm{d}\Omega\\
% &+\int_{\Gamma}v^{h}_{,n}M_n(T^{q}w)\mathrm{d}\Gamma+\left[v^{h}F_c(T^{q}w)\right]_{\Gamma_{c}}
% \end{aligned}
% \end{equation}
% In the meantime, by invoking a regular solution $w$ and eliminating the equivalent terms in Eq. \eqref{eq:35}, $\mathbf{A}^{*}(v^{h},w)$ and $\mathbf{F}^{*}(v^{h})$ can be recast as:
% \begin{equation}\label{eq:55}
% \bar{\mathbf{A}}^{*}(v^{h},u)=\int_{\Omega}^{*}B^{T}(v^{h})M(w)\mathrm{d}\Omega
% \end{equation}
% \begin{equation}\label{eq:56}
% \bar{\mathbf{F}}^{*}(v^{h})=-\int_{\Gamma}v^{h}V(w)\mathrm{d}\Gamma+\int_{\Omega}v^{h}q(w)\mathrm{d}\Omega+\int_{\Gamma}v^{h}_{,n}M_n(w)\mathrm{d}\Gamma+\left[v^{h}F_c(w)\right]_{\Gamma_{c}}
% \end{equation}
% It is noted that $\bar{\mathbf{F}}^{*}(v^{h})=\mathbf{B}^{*}(v^{h},w)$, $\forall v^{h}\in V^{h}$.
% Without consideration for numerical integration, we also naturally have the following identities:
% \begin{equation}\label{eq:57}
% \mathrm{A}(v^{h},w)=\mathrm{B}(v^{h},w),\quad \forall v^{h}\in V^{h}
% \end{equation}
% \begin{equation}\label{eq:58}
% \mathrm{A}(v^{h},T^{q}w)=\mathrm{B}(v^{h},T^{q}w),\quad \forall v^{h}\in V^{h}
% \end{equation}
% It is well known that the traditional integration, such as Gauss integration, based meshfree methods can not satisfy Eq. \eqref{eq:51}.
% However, the reproducing kernel smoothing gradient $\tilde{\Psi}_{I,ij}$ obtained by Eq. \eqref{eq:44} can naturally meets integration constraints \eqref{eq:37} and \eqref{eq:51}:
% \begin{equation}\label{eq:59}
% \begin{aligned}
% \mathbf{A}^{*}(v^{h},T^{q}w)&=\int_{\Omega}^{*}B^{T}(v^{h})M(T^{q}w)\mathrm{d}\Omega\\
% &=\alpha_{ij}^{T}\int_{\Omega}^{*}v^{h}_{,ij}p^{[q-2]T}\mathrm{d}\Omega\\
% &=\alpha_{ij}^{T}\sum_{I=1}^{n_p}v_I\int_{\Omega}^{*}\tilde{\Psi}_{I,ij}p^{[q-2]T}\mathrm{d}\Omega\\
% &=\alpha_{ij}^{T}\sum_{I=1}^{n_p}\sum_{C=1}^{NC}\underbrace{\left[\int_{\Omega_C}^{*}p^{[q-2]}p^{[q-2]T}\mathrm{d}\Omega\right]G_C^{-1}g_{ijI}^{C}}_{=\,G_C}v_I\\
% &=\alpha_{ij}^{T}\sum_{I=1}^{n_p}\sum_{C=1}^{NC}g_{ijI}^{C}v_I
% =\alpha_{ij}^{T}\sum_{I=1}^{n_p}g_{ijI}v_I\\
% &=\int_{\Gamma}v^{h}V(T^{q}w)\mathrm{d}\Gamma+\int_{\Omega}v^{h}q(T^{q}w)\mathrm{d}\Omega\\
% &\quad+\int_{\Gamma}v^{h}_{,n}M_n(T^{q}w)\mathrm{d}\Gamma
% +\sum_{m=1}^{n_c}\left[v^{h}F_c(T^{q}w)\right](c_m)\\
% &=\mathbf{B}^{*}(v^{h},T^{q}w)
% \end{aligned}
% \end{equation}
% In light of Eqs. \eqref{eq:51}, \eqref{eq:57} and \eqref{eq:58}, the orthogonality condition in Eq. \eqref{eq:50} can be rewritten as \cite{ref50}:
% \begin{equation}\label{eq:60}
% \begin{aligned}
% \mathrm{A}(v^{h},\tilde{e})
% &=\left(\mathbf{A}^{*}-\mathrm{A}\right)(v^{h},\tilde{e})+\mathbf{A}^{*}(v^{h},\tilde{e})\\
% &=\left(\mathbf{A}^{*}-\mathrm{A}\right)(v^{h},\tilde{e})+\mathbf{A}^{*}(v^{h},w)-\mathbf{A}^{*}(v^{h},w^{h})\\
% &=\left(\mathbf{A}^{*}-\mathrm{A}\right)(v^{h},\tilde{e})+\mathbf{A}^{*}(v^{h},w)-\mathbf{B}^{*}(v^{h},w)\\
% &=\left(\mathbf{A}^{*}-\mathrm{A}\right)(v^{h},\tilde{e})+\mathbf{A}^{*}(v^{h},w)-\mathbf{B}^{*}(v^{h},w)
% -\underbrace{\mathbf{A}^{*}(v^{h},T^{q}w)-\mathbf{B}^{*}(v^{h},T^{q}w)}_{=\,0}\\
% &=\left(\mathbf{A}^{*}-\mathrm{A}\right)(v^{h},\tilde{e})+\mathbf{A}^{*}\left(v^{h},w-T^{q}w\right)-\mathbf{B}^{*}\left(v^{h},w-T^{q}w\right)\\
% &=\left(\mathbf{A}^{*}-\mathrm{A}\right)(v^{h},\tilde{e})-\left(\mathbf{A}^{*}-\mathrm{A}\right)\left(v^{h},w-T^{q}w\right)\\
% &\quad+\left(\mathrm{B}-\mathrm{B}^{*}\right)\left(v^{h},w-T^{q}w\right)
% +\underbrace{\mathrm{A}(v^{h},w-T^{q}w)-\mathrm{B}(v^{h},w-T^{q}w)}_{=\,0}\\
% &\le\left|\left(\mathbf{A}^{*}-\mathrm{A}\right)(v^{h},\tilde{e})\right|
% +\left|\left(\mathbf{A}^{*}-\mathrm{A}\right)\left(v^{h},w-T^{q}w\right)\right|
% +\left|\left(\mathrm{B}-\mathrm{B}^{*}\right)\left(v^{h},w-T^{q}w\right)\right|
% \end{aligned}
% \end{equation}
% Under the premise of $(k-l)$th order of integration accuracy, following Cauchy--Schwarz inequality, the error bounds of three terms in Eq. \eqref{eq:60} can be stated as follows:
% \begin{equation}\label{eq:61}
% \left|\left(\mathbf{A}^{*}-\mathrm{A}\right)(v^{h},\tilde{e})\right|\le C_{\sigma}h\left|\mathrm{A}(v^{h},\tilde{e})\right|
% \end{equation}
% \begin{equation}\label{eq:62}
% \left|\left(\mathbf{A}^{*}-\mathrm{A}\right)\left(v^{h},w-T^{q}w\right)\right|\le C_{\sigma}h^{k-1}\left\|v^{h}\right\|_{H_2}\|w\|_{H_{k+1}}
% \end{equation}
% \begin{equation}\label{eq:63}
% \left|\left(\mathrm{B}-\mathrm{B}^{*}\right)\left(v^{h},w-T^{q}w\right)\right|\le C_{\sigma}h^{k-1}\left\|v^{h}\right\|_{H_2}\|w\|_{H_{k+1}}
% \end{equation}
% where $k=\min\{q,m\}$, $C_{\sigma}$ is a constant related to the numerical integration scheme.

% The proof of the error estimates of Eqs. \eqref{eq:61}--\eqref{eq:63} is given in the Appendix.

% Combining Eqs. \eqref{eq:60}--\eqref{eq:63}, the error estimate of orthogonality condition in Eq. \eqref{eq:50} can be given by:
% \begin{equation}\label{eq:64}
% \left|\mathrm{A}(v^{h},\tilde{e})\right|\le C_{\sigma}h^{k-1}\left\|v^{h}\right\|_{H_2}\|w\|_{H_{k+1}},\quad k=\min\{q,m\}
% \end{equation}
% It is noted that inequality of Eq. \eqref{eq:64} only can be ensured when numerical integration methods satisfy the $q$th order integration constraints.

% Subsequently, by introducing Eq. \eqref{eq:64}, the $H_2$ norm of error between final solution $\bar{u}^{h}$ and numerical solution $u^{h}$ of Kirchhoff--Love plate can be measured by:
% \begin{equation}\label{eq:65}
% \begin{aligned}
% \|\eta\|^{2}_{H_2}&\le(1/\alpha)\left|\mathrm{A}(\eta,\eta)\right|\\
% &=(1/\alpha)\left|\mathrm{A}(\eta,\tilde{e})-\underbrace{\mathrm{A}(\eta,e)}_{=\,0}\right|\\
% &=(1/\alpha)\left|\mathrm{A}(\eta,\tilde{e})\right|\\
% &\le C_{1q}h^{l-1}\|\eta\|_{H_2}\|w\|_{H_{l+1}},\quad l=\min\{q,m\}
% \end{aligned}
% \end{equation}
% where $\alpha$ and $C_{1q}$ both are positive constants unrelated to $h$.

% A combination of Eqs. \eqref{eq:49} and \eqref{eq:65} leads to:
% \begin{equation}\label{eq:66}
% \|\tilde{e}\|_{H_2}\le C_{1p}h^{k-1}\|u\|_{H_{k+1}}+C_{1q}h^{l-1}\|u\|_{H_{l+1}},\quad k=\min\{p,m\},\;l=\min\{q,m\}
% \end{equation}
% Similar to the accuracy analysis for second-order problems \cite{ref50}, it is noted that the error estimate for Galerkin meshfree methods of Kirchhoff--Love plate is mutually controlled by two parts.
% The first part comes from the standard interpolation error \eqref{eq:49}, in which the convergence rates is consistent with Galerkin isogeometric accuracy analysis \cite{ref50}; the second part comes from the loss of Galerkin orthogonality condition \eqref{eq:64}, caused by integration error of meshfree formulation.

% According to Eq. \eqref{eq:51}, the linear Lagrange interpolation ($T^{1}u\in\mathbb{P}^{1}$) will be eliminated in Eqs. \eqref{eq:53} and \eqref{eq:54}, thus the linear-order integration constraint ($q=1$) can be naturally satisfied no matter which integration scheme is adopted.
% Meanwhile, it is obvious that the traditional integration methods based meshfree formulation (such as Gauss integration) can hardly meet the higher-order integration constraints corresponding to the degree of basis functions ($p\ge2$, $q<p$).
% As a result, the corresponding $H_2$ error estimates for Galerkin meshfree formulations adopting different integration methods are listed in Table \ref{tab:1}.

% \begin{table}[pos=htbp]
% \centering
% \caption{Order of integration consistency and the corresponding convergence rate of $H_2$ error.}\label{tab:1}
% \begin{tabular}{@{}llcc@{}}
% \toprule
% \multirow{2}{*}{Order of integration consistency $q$} & \multirow{2}{*}{Integration scheme} & \multicolumn{2}{c}{Convergence rate of $H_2$ error}\\
% \cmidrule(lr){3-4}
% & & GI $(q=0)$ & RKGSI $(q=p)$\\
% \midrule
% $p\ge2$ & & 0 & $p-1$\\
% \bottomrule
% \end{tabular}
% \end{table}

% \subsection{$L_2$ norms of error for Galerkin meshfree formulation to Kirchhoff--Love plate}\label{sec:l2}

% In order to develop the accuracy analysis for Kirchhoff--Love plate, we further investigate the $H_1$ and $L_2$ error accounting for integration constraints by designing a series of auxiliary problems.

% Consider the following second-order scalar-valued problem:
% \begin{equation}\label{eq:67}
% \begin{cases}
% \Delta z=\tilde{e} & \text{in }\Omega\\[2pt]
% z=0,\quad z_{,n}=0 & \text{on }\Gamma
% \end{cases}
% \end{equation}
% where the source term is set as the displacement error $\tilde{e}$.

% According to the closed graph theorem, Eq. \eqref{eq:67} leads to the following inequality:
% \begin{equation}\label{eq:68}
% |z|_{H_1}\le C\|\tilde{e}\|_{L_2}
% \end{equation}
% Also, the Galerkin weak form of Eq. \eqref{eq:67} are given by:
% \begin{equation}\label{eq:69}
% \text{Find }z\in H_1(\Omega):\quad \mathrm{A}_1(v,z)=\mathrm{F}_1(v),\quad \forall v\in H_1(\Omega)
% \end{equation}
% Substituting $v=\tilde{e}$ into Eq. \eqref{eq:69}, $\mathrm{A}_2$ and $\mathrm{F}_2$ can be further expressed as:
% \begin{equation}\label{eq:70}
% \mathrm{A}_2(\tilde{e},g)=\int_{\Omega}\Delta\tilde{e}\,\Delta g\mathrm{d}\Omega,\quad
% \mathrm{F}_2(\tilde{e})=\int_{\Omega}\tilde{e}^{2}\mathrm{d}\Omega
% \end{equation}
% Introducing $T^{s<4}g\in\mathbb{P}^{s<4}$ to $\mathrm{A}_2(\tilde{e},g)$ yields the following identity:
% \begin{equation}\label{eq:71}
% \mathrm{A}_2\left(\tilde{e},T^{s<4}g\right)=\int_{\Omega}\tilde{e}\,\Delta\left(T^{s<4}g\right)\mathrm{d}\Omega=0
% \end{equation}
% where $s\le p$, $p$ is the order of meshfree interpolation for $w$.

% Thereafter, following Eq. \eqref{eq:71}, the following equality about $\|\tilde{e}\|^{2}_{L_2}$ holds:
% \begin{equation}\label{eq:72}
% \begin{aligned}
% \|\tilde{e}\|^{2}_{L_2}&=\mathrm{F}_2(\tilde{e})\\
% &=\mathrm{A}_2\left(w-w^{h},g\right)\\
% &=\mathrm{A}_2\left(w-w^{h},g\right)-\underbrace{\mathrm{A}_2\left(w-w^{h},T^{s<4}g\right)}_{=\,0}\\
% &=\mathrm{A}_2\left(w-w^{h},g-T^{s<4}g\right)
% \end{aligned}
% \end{equation}
% It is noted that, since the order selection $s$ of Lagrange interpolation is limited by the order of meshfree interpolation ($s\le p$), while the order of integration consistency $q$ will not affect the selection of $s$ in Eq. \eqref{eq:72}.
% Thus in order to differentiate the effect of different orders of meshfree interpolation, Eq. \eqref{eq:72} is divided into the following two parts:
% \begin{equation}\label{eq:73}
% \begin{aligned}
% \|\tilde{e}\|^{2}_{L_2}&=\mathrm{A}_2\left(w-w^{h},g-T^{s<4}g\right)\\
% &\le\left|\mathrm{A}_2\left(w-w^{h},g-T^{s<4}g\right)\right|+\left|\mathrm{A}_2\left(w^{h}-\bar{w}^{h},g-T^{s<4}g\right)\right|\\
% &=\left|\mathrm{A}_2\left(e,g-T^{s<4}g\right)\right|+\left|\mathrm{A}_2\left(\eta,g-T^{s<4}g\right)\right|
% \end{aligned}
% \end{equation}
% in which the error bound of second term is unrelated to the meshfree interpolation, thus $s=3$ is selected to determine the error estimate of $|\mathrm{A}_2(\eta,g-T^{s<4}g)|$:
% \begin{equation}\label{eq:74}
% \begin{aligned}
% \left|\mathrm{A}_2\left(\eta,g-T^{3}g\right)\right|&\le\gamma\|\eta\|_{H_2}\left\|g-T^{3}g\right\|_{H_2}\\
% &\le\gamma\|\eta\|_{H_2}C_{10}h^{2}|g|_{H_4}\\
% &\le C_{11}h^{2}\|\eta\|_{H_2}\left\|w-w^{h}\right\|_{L_2}
% \end{aligned}
% \end{equation}
% in which $C_{10}$ and $C_{11}$ both are constants.

% Regarding to the first term in Eq. \eqref{eq:73}, since $p$ constrains the selection of $s$, it can be further divided into two cases: $p=2$ and $p\ge3$.
% In order to get desirable error estimate for $\|\tilde{e}\|_{L_2}$, the interpolation order of $g$ can be respectively selected as different value in different cases.

% \noindent (1) $p=2$, $s=2$

% With continuity condition and Poincar\'e inequality, Eq. \eqref{eq:72} becomes:
% \begin{equation}\label{eq:75}
% \begin{aligned}
% \left|\mathrm{A}_2\left(\tilde{e},g-T^{2}g\right)\right|&\le\gamma\left\|w-w^{h}\right\|_{H_2}\left\|g-T^{2}g\right\|_{H_2}\\
% &\le\gamma\|\tilde{e}\|_{H_2}C_{20}h|g|_{H_3}\\
% &\le\gamma\|\tilde{e}\|_{H_2}C_{21}h|g|_{H_4}
% \end{aligned}
% \end{equation}
% where Poincar\'e inequality $\|g\|_{H_3}\le\alpha\|g\|_{H_4}$ is adopted, $C_{20}$ and $C_{21}$ both are constants.

% Introducing Eq. \eqref{eq:49} and Eq. \eqref{eq:65} into Eq. \eqref{eq:74} and Eq. \eqref{eq:75} respectively, then we have $L_2$ error estimate for Galerkin meshfree formulation with $p=2$:
% \begin{equation}\label{eq:76}
% \|\tilde{e}\|_{L_2}\le C_{2p}h^{k}|w|_{H_{k+1}}+C_{2q}h^{l+1}|w|_{H_{l+1}},\quad k=\min\{2,m\},\;l=\min\{q,m\}
% \end{equation}

% \noindent (2) $p\ge3$, $s=3$

% Following continuity condition, Eq. \eqref{eq:72} is further given by:
% \begin{equation}\label{eq:77}
% \begin{aligned}
% \left|\mathrm{A}_2\left(\tilde{e},g-T^{3}g\right)\right|&\le\gamma\left\|w-w^{h}\right\|_{H_2}\left\|g-T^{3}g\right\|_{H_2}\\
% &\le\gamma\left\|w-w^{h}\right\|_{H_2}C_{22}h^{2}|g|_{H_4}\\
% &\le C_{23}h^{2}\|\tilde{e}\|_{H_2}\left\|w-w^{h}\right\|_{L_2}
% \end{aligned}
% \end{equation}
% where $C_{22}$ and $C_{23}$ both are constants.

% Similarly, a combination of Eqs. \eqref{eq:74}, \eqref{eq:75} and \eqref{eq:77} leads to $L_2$ error estimate for Galerkin meshfree formulation with $p\ge3$:
% \begin{equation}\label{eq:78}
% \|\tilde{e}\|_{L_2}\le C_{3p}h^{k+1}|w|_{H_{k+1}}+C_{3q}h^{l+1}|w|_{H_{l+1}},\quad k=\min\{p,m\},\;l=\min\{q,m\}
% \end{equation}

% \subsection{$H_1$ norm of error for Galerkin meshfree formulation to Kirchhoff--Love plate}\label{sec:h1}

% Accordingly, in order to obtain $H_1$ error estimate for Galerkin meshfree formulation, we consider the following fourth-order auxiliary problem:
% \begin{equation}\label{eq:79}
% \begin{cases}
% \Delta z=\tilde{e} & \text{in }\Omega\\[2pt]
% z=0,\quad z_{,n}=0 & \text{on }\Gamma
% \end{cases}
% \end{equation}
% in which $\nabla w$ is a second-order tensor, the source term $\tilde{e}$ is defined as the gradient of displacement error accounting for integration error:
% \begin{equation}\label{eq:80}
% \tilde{e}=\nabla\left(w-\bar{w}^{h}\right)
% \end{equation}
% According to the closed-graph theorem, we have following relationship:
% \begin{equation}\label{eq:81}
% |z|_{H_1}\le C\|\tilde{e}\|_{L_2}
% \end{equation}
% where $C$ is a positive constant.

% By introducing the integration by parts, the Galerkin weak form of Eq. \eqref{eq:79} can be stated as:
% \begin{equation}\label{eq:82}
% \text{Find }z\in H_1(\Omega):\quad \mathrm{A}_1(v,z)=\mathrm{F}_1(v),\quad \forall v\in H_1(\Omega)
% \end{equation}
% where the bilinear form $\mathrm{A}_1$ and linear form $\mathrm{F}_1$ are defined as follows:
% \begin{equation}\label{eq:83}
% \mathrm{A}_1(v,z)=\int_{\Omega}\nabla v:\nabla z\mathrm{d}\Omega,\quad
% \mathrm{F}_1(v)=\int_{\Omega}v\cdot\tilde{e}\mathrm{d}\Omega
% \end{equation}
% Let the test function $v=\tilde{e}$, Eq. \eqref{eq:83} can be recast as:
% \begin{equation}\label{eq:84}
% \mathrm{A}_1(\tilde{e},z)=\int_{\Omega}\nabla\tilde{e}:\nabla z\mathrm{d}\Omega,\quad
% \mathrm{F}_1(\tilde{e})=\int_{\Omega}\tilde{e}\cdot\tilde{e}\mathrm{d}\Omega
% \end{equation}
% Invoking $s$th order Lagrange interpolation of $z$, since $T^{s<2}z\in\mathbb{P}^{s<2}$, the following identity holds:
% \begin{equation}\label{eq:85}
% \mathrm{A}_1\left(\tilde{e},T^{s<2}z\right)=\int_{\Omega}\tilde{e}\,\Delta\left(T^{s<2}z\right)\mathrm{d}\Omega=0
% \end{equation}
% where $s\le p-1$, $p$ is the order of meshfree interpolation for $w$.

% With Poincar\'e inequality, and Eq. \eqref{eq:85}, we have the following inequality:
% \begin{equation}\label{eq:86}
% \begin{aligned}
% \|\tilde{e}\|^{2}_{H_1}&=\|\tilde{e}\|^{2}_{L_2}+|\tilde{e}|^{2}_{H_1}\\
% &\le(\beta+1)|\tilde{e}|^{2}_{H_1}\\
% &=\alpha\mathrm{F}_1(\tilde{e})=\alpha\mathrm{A}_1(\tilde{e},z)\\
% &=\alpha\left|\mathrm{A}_1(\tilde{e},z)-\underbrace{\mathrm{A}_1\left(\tilde{e},T^{s<2}z\right)}_{=\,0}\right|\\
% &=\alpha\left|\mathrm{A}_1\left(\tilde{e},z-T^{s<2}z\right)\right|\\
% &\le\alpha\left|\mathrm{A}_1\left(e,z-T^{s<2}z\right)\right|+\alpha\left|\mathrm{A}_1\left(\eta,z-T^{s<2}z\right)\right|
% \end{aligned}
% \end{equation}
% where $\alpha=\beta+1$, $\beta$ is a positive constant.

% It is noted that $s=1$ is the only choice for evaluating the error bounds of both first and second terms, thus there is no need to divide $\|\tilde{e}\|^{2}_{H_1}$ into two parts.
% Subsequently, by combining Eqs. \eqref{eq:81} and \eqref{eq:86}, with continuity condition and Poincar\'e inequality, the $H_1$ error estimate for Galerkin meshfree formulation to Kirchhoff--Love plate can be given by:
% \begin{equation}\label{eq:87}
% \begin{aligned}
% |\tilde{e}|^{2}_{H_1}&\le\alpha\left|\mathrm{A}_1\left(\tilde{e},z-T^{1}z\right)\right|\\
% &\le\alpha\gamma\|\tilde{e}\|_{H_1}\left\|z-T^{1}z\right\|_{H_1}\\
% &\le\gamma\left\|w-w^{h}\right\|_{H_2}C_{30}h|z|_{H_2}\\
% &\le\gamma C_{31}h\left\|w-w^{h}\right\|_{H_2}\|\tilde{e}\|_{L_2}\\
% &\le C_{32}h\left\|w-w^{h}\right\|_{H_2}\left\|w-w^{h}\right\|_{H_1}
% \end{aligned}
% \end{equation}
% where $C_{30}$, $C_{31}$ and $C_{32}$ both are constants.

% By further introducing Eq. \eqref{eq:80}, Eq. \eqref{eq:87} becomes:
% \begin{equation}\label{eq:88}
% \|\tilde{e}\|_{H_1}\le C_{4p}h^{k}|u|_{H_{k+1}}+C_{4q}h^{l}|u|_{H_{l+1}},\quad k=\min\{p,m\},\;l=\min\{q,m\}
% \end{equation}
% Now we have obtained the $H_2$ error estimate shown in Eq. \eqref{eq:49}, $L_2$ error estimate shown in Eqs. \eqref{eq:76} and \eqref{eq:78}, and $H_1$ error estimate shown in Eq. \eqref{eq:88} for Galerkin meshfree formulation to Kirchhoff--Love plate, i.e., an explicit accuracy assessment accounting for the effect of integration constraints of Kirchhoff--Love plate has been investigated with $H_1$, $H_2$ and $L_2$ norms.
% The corresponding theoretical convergence rates for different integration schemes and interpolation order are listed in Table \ref{tab:2}.

% \begin{table}[pos=htbp]
% \centering
% \caption{Theoretical convergence rates of the $H_2$, $H_1$ and $L_2$ errors for GI and RKGSI.}\label{tab:2}
% \begin{tabular}{@{}llcc@{}}
% \toprule
% \multicolumn{2}{@{}l}{} & \multicolumn{2}{c}{Integration scheme}\\
% \cmidrule(lr){3-4}
% \multicolumn{2}{@{}l}{} & GI & RKGSI\\
% \midrule
% \multicolumn{2}{@{}l}{Order of integration consistency} & $q=1$ & $q=p$\\
% \midrule
% Convergence rate of $H_2$ error & $p\ge2$ & 0 & $p-1$\\
% Convergence rate of $H_1$ error & $p\ge2$ & 1 & $p$\\
% Convergence rate of $L_2$ error & $p=2$ & 2 & $p$\\
%                                  & $p\ge3$ & 2 & $p+1$\\
% \bottomrule
% \end{tabular}
% \end{table}

% It is noted in Table \ref{tab:2} that the integration schemes play a crucial role in Galerkin meshfree formulation, only when they satisfy $p$th order of integration constraints, can they achieve the optimal convergence performance.
% It is also noted that, the convergence rates $(p+1)$ for $L_2$ and $p$ for $H_1$ error can only be attained when the meshfree interpolation order is greater than two.
